\documentclass[12pt]{article}
\usepackage[utf8]{inputenc}
\usepackage[T1]{fontenc}
\usepackage[spanish]{babel}
\usepackage{geometry}
\geometry{margin=2cm}
\usepackage{parskip}
\usepackage{pdfpages}
\usepackage{xcolor}
\usepackage{graphicx}
\usepackage{float}
\usepackage{minted}
\usepackage{enumitem}
\usepackage{listings}
\lstset{
  backgroundcolor=\color{gray!10},
  basicstyle=\ttfamily\small,
  frame=single,
  breaklines=true,
  keywordstyle=\color{blue}
}

\begin{document}
\begin{titlepage}
\begin{center}
    % Logos
    \begin{minipage}{0.3\textwidth}
        \includegraphics[width=0.8\linewidth]{UNAM_Logo(1).png}
    \end{minipage}\hfill
    \begin{minipage}{0.4\textwidth}
        \begin{center}
            \textbf{\Large Universidad Nacional Autónoma de México}\\[0.2cm]
            \textbf{Facultad de Ciencias, 2027-1}\\[0.2cm]
        \textbf{Fundamentos de Bases de Datos}
        \end{center}
    \end{minipage}\hfill
    \begin{minipage}{0.3\textwidth}
        \includegraphics[width=0.8\linewidth]{Logo_FC(1).png}
    \end{minipage}
    
    \vspace{1cm}
    \textbf{\Large READ ME \\ 
    Los Fundamentados Basados Informados\\ 
    (LFBI)}\\
    \vspace{1cm}
    \large{Integrantes:\\
    Diego Alexander Cisneros Tenorio (322251050)\\
    José Manuel Nuñez Ruiz (322112072)\\
    Andrea Lucia Olmos Ortiz (322149977)\\
    Alexa Reynoso Ascencio (322179457)\\
    Saúl Vega Navas (322088267}\\
    \vspace{1cm}
    \textbf{Profesor: }\\
    Gerardo Avilés Rosas\\
    \vspace{0.5cm}
    \textbf{Ayudantes de Teoría: }\\
    Oscar Daniel Flores Linares\\
    Claudia Itzel Gutiérrez Sánchez \\
    \vspace{0.5cm}
    \textbf{Ayudantes de Laboratorio: }\\
    Ricardo Badillo Macias \\
    Rocío Aylin Huerta González \\
    \vspace{2cm}    
    \textbf{Fecha de entrega: } 07 de Septiembre de 2026
    \vfill
\end{center}
\end{titlepage}
En esta entrega se tienen dos carpetas:
\begin{itemize}
    \item Docs: Con un archivo Practica02.pdf que contiene el análisis de requerimientos para el caso de uso del "PuellaGame"
    \item SRC: Con el programa correspondiente al caso de uso "PuellaGame" y una carpeta Doc con la documentación correspondiente al programa realizado.
\end{itemize}
El programa fue implementado en Java y en la carpeta SRC se incluyen los archivos compilados (.class) correspondientes a cada programa, de este modo para ejecutarlo basta con abrir la terminal, ubicarse en la carpeta SRC y ejecutar:
\begin{itemize}
    \item \texttt{java main.java}
\end{itemize}
Para el programa se agregaron algunos premios, clientes y sucursales sugeridos para probar el funcionamiento, no obstante el sistema no está limitado a los mismos y es posible eliminarlos y guardar nueva información. A continuación se enlistan los ejemplares que ya vienen en el sistema listos para ser consultados.\\

\subsection*{Archivo de Clientes}
\textbf{Orden} \texttt{idCliente, curp, nombre, apellidoPaterno, apellidoMaterno, fechaNacimiento, edad, sexo, correos, telefonos}
\begin{itemize}
    \item \textbf{01}, REAA030515MDFRSC01, Alexa, Reynoso, Ascencio, 15/05/2003, 23, F, alexa@example.com, 5512345678
    \item \textbf{02}, CIDA020820HDFSTN02, Diego Alexander, Cisneros, Tenorio, 20/08/2002, 24, M, diego@example.com, 5587654321
    \item \textbf{03}, OOAL031110MDFRRT03, Andrea Lucia, Olmos, Ortiz, 10/11/2003, 22, F, lucia@example.com, 5523456789
    \item \textbf{04}, VENS010305HDFGVN04, Saul, Vega, Navas, 05/03/2001, 25, M, saul@example.com, 5534567890
    \item \textbf{05}, NURJ030130HDFNZR05, José Manuel, Nuñez, Ruiz, 30/01/2003, 23, M, jose@example.com, 5545678901
\end{itemize}

\subsection*{Archivo de Premios}
\textbf{Orden} \texttt{idPremio, Nombre, Categoría, Rango de Edad, Valor Aproximado, Puntos de Canjeo, Cantidad Disponible}
\begin{itemize}
    \item \textbf{1001}, Elefante de peluche, Medio, Infantil, 1000.0, 2000, 5
    \item \textbf{1002}, Patineta azul, Medio, Juvenil, 1750.0, 3500, 3
    \item \textbf{1003}, Dulces, Bajo, Infantil, 10.0, 20, 80
    \item \textbf{1004}, Set de Legos , Medio, Infantil, 1350.0, 2700, 14
    \item \textbf{1005}, Audifonos rojos, Alto, Juvenil, 3500.0, 7000, 78
    \item \textbf{1006}, Playera, Alto, Juvenil, 2000.0, 4000, 80
    \item \textbf{1007}, Telefono, Alto, +Adulto, 5000.0, 10000, 5
    \item \textbf{1008}, Conejo de Peluche, Alto, Infantil, 200.0, 400, 60
\end{itemize}

\subsection*{Archivo de Sucursales}
\textbf{Orden} \texttt{idSucursal, nombre, calle, numeroInt, numeroExt, colonia, estado, telefono, horarioApertura, horarioCierre}
\begin{itemize}
    \item \textbf{1}, PuellaGame CU, Invetigación Cientifica, S/N, 04510, Coyoacán, Ciudad de México, 5583748932, 9:00, 20:00
    \item \textbf{3}, PuellaGame Aragon, Av. Carlos Hank Gonzáles, S/N, 120, Ecatepect de Morelos, Estado de México, 5513239904, 9:00, 10:00
    \item \textbf{5}, PuellaGame Lindavista, Av Montevideo, S/N, 363, Gustavo A. Madero, Ciudad de México, 5557263100, 11:00, 8:30
    \item \textbf{6}, PuellaGame Pilgrim, Puella, 25, 432, Liverpool, Guadalajara, 5599382819, 7:00, 5:00
\end{itemize}


\end{document}
